\documentclass[11pt]{article}
\usepackage[a4paper,margin=24mm]{geometry}
\usepackage[T1]{fontenc}
\usepackage{lmodern,amsmath,amssymb,amsthm,mathtools,booktabs,microtype}
\usepackage[colorlinks=true,linkcolor=blue,urlcolor=blue,citecolor=blue]{hyperref}
\newcommand{\tr}{\tau}
\newcommand{\A}{\mathcal A}
\newtheorem{proposition}{Proposition}
\title{A preparation dictionary for the selected state\\\large Modular spectra and energy-resolved balance}
\author{Alan Day \and Chengheng\thanks{AI research assistant (OpenAI).}}
\date{8 November 2026\quad Local research draft}
\begin{document}\maketitle
\begin{abstract}
The complete selected bilocal mean and covariance call for a
distinction between preparation and the generator of static-patch
time. We apply standard tracial modular theory to the already
established selected density. Its modular flow has an explicit
q-Mehler spectral representation after an orthogonal rotation of
the two q-Gaussian fields. We derive a nonzero stationary part in
the modular correlation of the microscopic Hamiltonian, spectral
sum rules, and a finite-window test that rules out every constant
Hamiltonian temperature even on the stationary charged restriction.
These are model-level preparation data for a subsequent bulk
dictionary; no spatial generator is supplied by taking a logarithm
of the density.
\end{abstract}

\section{The physical comparison and its conventions}
Miyashita--Sekino--Susskind identify the microscopic Hamiltonian
with static-patch boost time, while distinguishing its Boltzmann
temperature from the Gibbons--Hawking temperature
\cite[Sec.~2, (2.4)--(2.5)]{MSS26}.
Their flat-space construction uses the maximally mixed microscopic
background, whereas the bulk Rindler vacuum has dimensionless
temperature $1/(2\pi)$ \cite[Secs.~1, 6]{MSS25}.
The singlet mass comparison is still a quantitative target
\cite[Sec.~4.3]{MSS26}. These statements do not assign a new
temperature to every prepared microscopic state.

Our model is the same-disorder overlapping complex Gaussian
$p$-body CAR system, at half filling. Exact finite-$N$ pinching
precedes $p^2/N\to\lambda>0$ double scaling. Use
\[
 q=e^{-4\lambda},\quad \zeta=\sqrt q,\quad C=J_s^2/\lambda,
 \quad h=H/\sqrt C,\quad W=\sqrt\lambda U,\quad\hbar=1.
\]
The fields $h,U$ are orthogonal unit-variance q-Gaussians in the
faithful tracial limiting algebra. Let $\A$ be the controlled
covariant algebra containing them and the charged curvature $X$.
The preceding selected-state and higher-descendant results give
\begin{equation}\label{density}
 \sigma(A)=\tr(\rho A),\qquad
 \rho=1-\zeta(1+\lambda)+\zeta(h-W)^2.
\end{equation}
All algebraic results below concern this limiting algebra. We do
not infer microscopic modular convergence from polynomial moment
convergence alone.

Physical Hamiltonian time is $\alpha_T(A)=e^{iTh}Ae^{-iTh}$,
with $T=\sqrt C\,t_{\rm phys}$. The previous string variable is
$u=2\sqrt\lambda T$; the earlier candidate Rindler matching used
$\eta=2u$. The modular parameter introduced below is a separate
dimensionless variable. Lin--Stanford's bra/ket maps
\cite[(158)--(159)]{LS} are kept as preparation-dependent maps;
they are not a uniform rescaling of all these clocks.

\newpage
\section{A computable modular flow}
Use the already established rotation
\begin{equation}\label{rotation}
 v=\frac{h-\sqrt\lambda U}{\sqrt{1+\lambda}},\qquad
 z=\frac{\sqrt\lambda h+U}{\sqrt{1+\lambda}},\qquad
 h=\frac{v+\sqrt\lambda z}{\sqrt{1+\lambda}}.
\end{equation}
Orthogonal second quantization preserves the q-Gaussian trace
\cite[Theorem 2.11]{BKS}; hence $v,z$ have the same joint law
as $h,U$. Set
\begin{equation}\label{ab}
 a=1-\zeta(1+\lambda),\quad b=\zeta(1+\lambda),\quad
 R=\frac2{\sqrt{1-q}},\quad p(x)=a+bx^2.
\end{equation}
Then $\rho=p(v)$, $a+b=1$ and $a>0$, since
$e^{2\lambda}>1+\lambda$. The known exact spectral bounds are
$a\le\rho\le a+bR^2$. In particular the state is faithful.

On the tracial standard Hilbert space $L^2(\A,\tr)$, the
normalized cyclic separating vector is $\Omega_\sigma=\rho^{1/2}$.
For left and right multiplication $L,R$, standard modular theory
\cite[Sec.~4.2]{Witten} gives
\[
 \Delta_\sigma=L_\rho R_{\rho^{-1}},\qquad
 \Delta_\sigma^{is}A\Delta_\sigma^{-is}=\rho^{is}A\rho^{-is}.
\]
We choose the opposite sign to use ordinary positive-temperature
Heisenberg conventions:
\begin{equation}\label{flow}
 k(x)=-\log p(x),\quad K=k(v),\qquad
 \beta_s(A)=e^{isK}Ae^{-isK}=\rho^{-is}A\rho^{is}.
\end{equation}
Then $\sigma$ is a $\beta=1$ KMS state for $\beta_s$.
The implementing generator on $L^2$ is $L_K-R_K$, not $L_K$
alone. This follows directly by moving $\rho$ around the trace;
no finite-dimensional matrix trace is assumed.

The state also has the exact preparation forms
\begin{equation}\label{prep}
 \sigma(A)=a\tr(A)+b\tr(vAv)
           =\langle\rho^{1/2},A\rho^{1/2}\rangle_2.
\end{equation}
These are two descriptions of the same established density.
The rotation, positivity and nonstationarity were already proved
in \cite{selected}; the new calculations below are its modular
response and an energy-resolved balance test.

In particular the inherited identity
\[
 \|[h,\rho]\|_2^2
 =2q\lambda(1+\lambda)(1+q)(1-q^2)>0
\]
shows that no nonzero constant rescaling of $\alpha_T$ equals
\eqref{flow} on the full algebra: $\beta_s$ preserves $\sigma$,
whereas $\alpha_T$ does not. The modular clock is available for
computing preparation response, not as a replacement boost.

\newpage
\section{Exact modular spectral kernel}
Let $d\mu_q$ be the unit-variance q-Hermite measure and let
$K_r(x,y)$ denote its Mehler kernel with eigenvalues $r^n$.
The one-$z$ bimodule measure relative to $v$ is
\[
 d\nu(x,y)=K_q(x,y)d\mu_q(x)d\mu_q(y).
\]
For $\mathcal W(s)=\sigma(\beta_s(z)z)$, functional calculus
and the two-$z$ trace formula give
\begin{equation}\label{kernel}
 \boxed{\mathcal W(s)=\int p(x)e^{is[k(x)-k(y)]}\,d\nu(x,y).}
\end{equation}
Here $x$ is the state-weighted energy of $v$, not the original
Hamiltonian energy $h$. The positive spectral measure, using
$\mathcal W(s)=\int e^{-i\omega s}dS(\omega)$, is the pushforward
of $p(x)d\nu$ under
\[
 \omega=k(y)-k(x)=\log\frac{p(x)}{p(y)}.
\]
It is supported in $[-\Omega,\Omega]$, with
\begin{equation}\label{band}
 \Omega=\log\frac{a+bR^2}{a}.
\end{equation}
Exchanging $x,y$ proves the measure identity
$dS(-\omega)=e^{-\omega}dS(\omega)$, equivalently
\begin{equation}\label{kms}
 \mathcal W(s-i)=\mathcal W(-s)\qquad(s\in\mathbb R).
\end{equation}
At fixed $q$, the kernel is entire in $s$ because $K$ is bounded.
This KMS statement fixes the sign convention in \eqref{flow}.

Some useful exact sum rules are
\begin{align}
 \mathcal W(0)&=1,\\
 i\mathcal W'(0)&=\frac12\int[p(x)-p(y)]
                  \log\frac{p(x)}{p(y)}d\nu>0,\label{first}\\
 -\mathcal W''(0)&=\int p(x)\left[\log\frac{p(x)}{p(y)}\right]^2d\nu,\\
 D(\sigma\|\tr)&=\int p(x)\log p(x)d\mu_q(x).
\end{align}
The entropy is relative to the normalized trace, not an absolute
von Neumann entropy of a continuum wedge. Strict positivity in
\eqref{first} uses the strictly positive kernel and nonconstant $p$.

The microscopic Hamiltonian has the modular correlation
\begin{equation}\label{hcor}
 \sigma(\beta_s(h)h)=
 \frac{1+b(1+q)}{1+\lambda}
       +\frac{\lambda}{1+\lambda}\mathcal W(s).
\end{equation}
Cross terms vanish by odd $z$ parity, and
$\sigma(v^2)=1+b(1+q)$. The $\mathcal W$ spectrum has no atoms:
on each half of $[-R,R]$, $k$ is strictly monotone with nonzero
derivative away from zero. Its pushforward of the integrable
two-variable density is absolutely continuous; the isolated
critical point carries no mass. Taking the difference of the
two transformed variables preserves absolute continuity.
Riemann--Lebesgue therefore gives $\mathcal W(s)\to0$.
Thus \eqref{hcor} has the explicit nonzero plateau shown there.
This is a modular-time result after fixed-$\lambda$ double scaling.

\newpage
\section{A finite-window test of Hamiltonian temperature}
The previously proved charged restriction on $W^*(h,X)$ is
\[
 \sigma|_{h,X}=\tau_f,\qquad
 f(h)=1-\zeta+\zeta h^2,
\]
where $X$ is the normalized centered charged curvature with
Hamiltonian crossing $q^2$. This restriction is stationary under
$\alpha_T$ even though the full state is not.
Its unfiltered $X$ autocorrelation is real and even. Consequently
that one autocorrelation alone has symmetric frequency weight
and cannot establish tracial equilibrium on the whole algebra.

Take a bounded energy-filtered probe
$A_{IJ}=P_I(h)XP_J(h)$. The exact same-disorder two-energy
measure is $d\nu_X=K_{q^2}(x,y)d\mu_q(x)d\mu_q(y)$, and
\begin{align}\label{rates}
 F_{IJ}:=\tau_f(A_{IJ}^\dagger A_{IJ})
  &=\int_{I\times J}f(y)d\nu_X,\\
 R_{IJ}:=\tau_f(A_{IJ}A_{IJ}^\dagger)
  &=\int_{I\times J}f(x)d\nu_X.
\end{align}
These are positive probe norms (or transition weights), not
rates obtained from a new interaction model.

\begin{proposition}
The restriction $\tau_f$ on $W^*(h,X)$ is not a KMS state for
$\alpha_T$ at any constant real inverse temperature $\beta$.
\end{proposition}
\begin{proof}
The bounded inner dynamics makes every bounded probe entire
analytic. Its KMS identity requires energy-resolved reverse to
forward ratio $e^{-\beta(x-y)}$. Equivalently,
\[
 R_{IJ}=\int_{I\times J}e^{-\beta(x-y)}f(y)d\nu_X.
\]
Choose $I_+=[.8,1.2]$, $I_-=[-1.2,-.8]$, $J=[-.2,.2]$.
All intervals lie inside the support for every $0<q<1$.
The positive Mehler kernel and monotonicity of $f$ in $|x|$ give
\begin{equation}\label{ratio}
 \frac{R_{I_+J}}{F_{I_+J}}
 =\frac{R_{I_-J}}{F_{I_-J}}
 \ge\frac{1-\zeta+.64\zeta}{1-\zeta+.04\zeta}>1.
\end{equation}
The equality follows from simultaneous reflection. In the
positive window $x-y>0$, so this inequality would force
$\beta<0$. In the negative window it would force $\beta>0$.
The value $\beta=0$ would give ratio one in both. No constant
temperature is consistent with the two probes.
\end{proof}

At the density level the same obstruction is
$f(x)/f(y)\ne e^{-\beta(x-y)}$ for any fixed $\beta$ on the
full support: evenness would force $\beta=0$, then $f$ would
have to be constant. The restricted state's own modular
Hamiltonian is instead $K_f=-\log f(h)$.
This distinguishes physical stationarity from thermal balance,
and uses actual probes in the charged algebra rather than an
assumed temperature inferred from a fit to one correlator.

\newpage
\section{Executed checks and error scope}
We evaluated \eqref{kernel} at $q=1/4,1/2$, with
$\lambda=-\log q/4$, using two independent representations:
the q-Hermite Jacobi matrix and the angular q-Pochhammer density
and Mehler product. Each used orders 128 and 192.
The 22 comparisons cover real modular times $0,.5,1,2,4$,
mid-strip times $.5-i/2,1-i/2,2-i/2$, and lower-boundary times
$.5-i,1-i,2-i$. All complex values are saved.

The maximum cross-representation absolute difference is
$2.34\cdot10^{-14}$, and the largest within-method refinement
is $1.88\cdot10^{-12}$. The maximum cross relative difference
is $2.76\cdot10^{-14}$; the largest phase-cancellation factor
is $72.92$. Neither comparison is an interval certificate.
KMS boundary exchanges are sign/implementation checks, not an
independent proof of modular theory. Sum rules were compared
between the two representations as well.

\begin{center}
\begin{tabular}{cccc}
\toprule
$q$&$\mathcal W(1)$&$-\mathcal W''(0)$&$R_{I_+J}/F_{I_+J}$\\
\midrule
$1/4$&$.55776147-.30164823i$&$1.09084006$&$1.97039949$\\
$1/2$&$.40881713-.31551359i$&$1.70797502$&$3.28486964$\\
\bottomrule
\end{tabular}
\end{center}
Four finite-window calculations (both signs and both $q$) used
32/64-point interval quadrature. Their ratios exceed the analytic
lower bounds $1.57692308$ and $2.32096452$ respectively. The
analytic incompatibility in \eqref{ratio} does not depend on
these numerical values.

\paragraph{A separately controlled product tail.}
Truncate the density products at $n=P$ and the kernel product
at $n=P-1$, as in the code. For a general Mehler parameter $r$,
the absolute logarithmic error of the joint density is bounded by
\begin{equation}\label{producttail}
 \delta_P=\frac1{1-q}\left[
 \frac{6q^{P+1}}{1-q^{P+1}}+
 \frac{r^2q^P}{1-r^2q^P}+
 \frac{4rq^P}{1-rq^P}\right].
\end{equation}
This follows termwise from
$|\log(1-z)|\le |z|/(1-|z|)$, including each conjugate factor.
For real $s$, the finite-product continuous integral differs
from \eqref{kernel} by at most $e^{\delta_P}-1$; its positive
reference weight $p(x)d\nu$ has mass one. For
$s=t-i\theta$, $0\le\theta\le1$, the absolute weight is
$p(x)^{1-\theta}p(y)^\theta d\nu$, whose mass is at most one
by weighted arithmetic--geometric mean. The same bound holds.
For the second sum rule multiply by $\Omega^2$.
We used $P=192$. This bound excludes angular quadrature,
Jacobi truncation and floating-point error.

All computations are fixed-$q$ diagnostics. The logarithmic
flow and its bounds are not uniform as $\lambda\downarrow0$;
no small-$\lambda$ modular limit is asserted.

\newpage
\section{What this changes in the bulk construction}
The constructive output is a definite preparation vector
$\rho^{1/2}$ and its exact modular spectral response. It allows
us to keep three objects separate in the next comparison:
the trace background used in the DSSYK proposal, the prepared
state used in our bilocal calculation, and the generator of
candidate bulk boost time. A stationary charged restriction does
not remove the need to match the preparation: its energy-filtered
balance test still differs from a constant-temperature Gibbs law.

For a vacuum wedge in relativistic QFT, the modular/boost
identification uses the wedge algebra and the vacuum, with the
fixed $2\pi$ factor \cite[Sec.~5]{Witten}. Applying this statement
to an arbitrary selected state and the entire microscopic algebra
would change both inputs. Here $K=k(v)$ is explicitly known,
but taking $K$ to be the geometrical boost is not the dictionary.

The next constructive target is the action of a candidate null
translation on the \emph{prepared} cyclic states, with the boost
still implemented by the original Hamiltonian time. The earlier
translation note \cite{translations} supplies a check:
two-point boost kernels alone admit arbitrary masses, while
reverse-null response probes mass moments. We should calculate
generator matrix elements with the bra/ket preparation kept
explicit, verify the mixed boost/translation Ward identities,
and only then compare the mass operator to the existing
't Hooft target. This stage supplies the state and balance data
for that test; the null generator, selected four-bilocal dynamics
and finite-$G$ dictionary remain to be constructed.

\begin{thebibliography}{9}\small
\bibitem{MSS26} S. Miyashita, Y. Sekino and L. Susskind,
\emph{Holograms and Standard Models},
\href{https://arxiv.org/abs/2607.05678v2}{arXiv:2607.05678v2},
Secs.~2 and 4.3.
\bibitem{MSS25} S. Miyashita, Y. Sekino and L. Susskind,
\emph{DSSYK at Infinite Temperature: The Flat-Space Limit and
the 't Hooft Model},
\href{https://arxiv.org/abs/2506.18054v3}{arXiv:2506.18054v3},
Secs.~1 and 6. The earlier phase-boundary location is corrected
in L. Susskind, \href{https://arxiv.org/abs/2511.10907}{arXiv:2511.10907}.
\bibitem{LS} H. W. Lin and D. Stanford,
\emph{A symmetry algebra in double-scaled SYK},
\href{https://arxiv.org/abs/2307.15725}{arXiv:2307.15725},
Sec.~4.2, (153)--(164).
\bibitem{BKS} M. Bo\.{z}ejko, B. K\"ummerer and R. Speicher,
\emph{q-Gaussian processes: non-commutative and classical aspects},
Commun. Math. Phys. 185 (1997), 129--154,
\href{https://arxiv.org/abs/funct-an/9604010}{arXiv:funct-an/9604010}.
\bibitem{Witten} E. Witten,
\emph{Notes on Some Entanglement Properties of Quantum Field Theory},
\href{https://arxiv.org/abs/1803.04993v6}{arXiv:1803.04993v6},
Secs.~4.1--4.2 and 5.
\bibitem{selected} A. Day and Chengheng,
\emph{Single selected Slater preparation}, local stage
\texttt{susskind-slater-selected-state}, 8 Nov. 2026,
and \texttt{susskind-slater-higher-wick} for the extended tower.
\bibitem{covariance} A. Day and Chengheng,
\emph{Preparation memory in the complete bilocal covariance},
local stage \texttt{susskind-selected-bilocal-covariance},
8 Nov. 2026.
\bibitem{translations} A. Day and Chengheng,
\emph{From boost kernels to candidate translations}, local stage
\texttt{susskind-translations}, 6 Nov. 2026.
\end{thebibliography}
\end{document}
